\scanpage{6}
Sauf erreur, on est dans les conditions qu'il faut pour que la \deleted{polyèdre comb.}\inserted{carte comb.} régulière définie par la $n$-carte $R\simeq R_{p_*}$ soit défini par une \uline{géométrie d'incidence} ayant \deleted{les} $D_i$ comme ens des facettes de dim $i$ (donc soit un polyèdre comb. régulier \inserted{$\Pi$}).

Une \uline{réalisation géom.} \inserted{\uline{affine régulière}} de $\Pi$ \deleted{est} sur un \deleted{schéma} anneau comm. $k$ (disons) est la donnée

a) d'un torseur (espace $k$-affine) \inserted{$E$} sous un $k$-Module loc libre de rang $n+1$

b) d'applications
\[
\begin{aligned}
\unsure{\varphi|\phi}_i:D_i\longrightarrow{}&\text{ens des sous-esp. affines}\\
&\text{de dim }i\text{ dans }E
\end{aligned}
\]
\[
(0\leq i\leq n)
\]
respectant les relations d'incidences,

\deleted{c) d'un hom. de groupes}\inserted{avec condition supplémentaire}

\unsure{par|$\rho$} \deleted{\unsure{par}}
\diagram{\begin{tikzcd}[ampersand replacement=\&]
\deleted{\operatorname{Aut}(\Pi)} \arrow[r,phantom,"\deleted{\longrightarrow}" description] \& \deleted{\operatorname{Autaff}(E)}
\end{tikzcd}}
\deleted{\unsure{tel}} \deleted{compatible avec les}

c)
\[
\left\{\begin{aligned}
&\forall g\in\operatorname{Aut}(\Pi),\quad\exists!\ \bar{g}\in\operatorname{Aff}(E)\quad\text{tel que}\\
&\bar{g}\circ\varphi_i=\varphi_i\circ g_{D_i}\qquad\forall\ 0\leq i\leq n.
\end{aligned}\right.
\]
--- et cette unicité est \unsure{vrai|vraie} fibre par fibre \dots
% archive sheet number: 5
